\section{Unary operations with role switching}\label{app:ambiguity}
Throughout this appendix, terms use the full signature
$\mathbf A_I^{\rm sym}$.

\subsection{Three roles: the exact unary normal form}

Now let $I=\{a,b,c\}$. A partial permutation $p$ is an injective partial
map from $I$ to itself. Put
\[
 C_p=I\setminus\operatorname{ran}(p),\qquad
 F_p(S)=C_p\cup p[S\cap\operatorname{dom}(p)].
\]

\begin{theorem}[All unary term operations]\label{sym:thm:unary}
Every unary term function is either a constant or has the following
unique normal form:
\begin{equation}\label{sym:eq:unary}
 f(S)=\begin{cases}D_i,&S=\{i\},\\F_p(S),&|S|\geq2,\end{cases}
\end{equation}
where $p$ is a partial permutation and the nonempty rows satisfy
\[
 \begin{array}{ll}
 \{p(i)\}\subseteq D_i\subseteq C_p\cup\{p(i)\},
       &i\in\operatorname{dom}(p),\\
 \varnothing\ne D_i\subseteq C_p,
       &i\notin\operatorname{dom}(p).
 \end{array}
\]
Every such normal form is realized by a term. The only constant also
in this family is the constant $I$.
\end{theorem}

We prove the upper bound, uniqueness, and synthesis separately. The finite checks are described in \cref{sec:verification}.

\subsubsection{Upper bound and mandatory images}

A unary term is either constant or admits a live provenance
$(C,P)$ with nonempty $P\subseteq\operatorname{Sym}(I)$ such that
\[
 f(S)=C\cup\bigcup_{\pi\in P}\pi S\quad(|S|\geq2),\qquad
 f(\{i\})\subseteq C\cup\{\pi(i):\pi\in P\}.
\]
This follows by the same induction as Lemma~\ref{sym:lem:provenance}, keeping
globally constant children separate. A constant owner singleton absorbs
globally; otherwise a live child remains live. In particular, a live
function preserves $I$ and maps ambiguous inputs to ambiguous outputs.

The induction also gives the following mandatory-image fact:
\begin{equation}\label{sym:eq:mandatory}
 k\notin C\text{ and }\pi(i)=k\text{ for every }\pi\in P
       \quad\Longrightarrow\quad k\in f(\{i\}).
\end{equation}
At an owned node, every live child satisfying this hypothesis contains
$k$ at input $\{i\}$. If the owner differs from $k$, such a child cannot
be its singleton; a constant child cannot be that owner singleton
either, since the whole expression would have been globally constant.
Thus the node preserves $k$. If the owner is $k$, membership of $k$ is
disjunction and is preserved directly. Role switching transports this
argument; the variable supplies its base case.

For an output role $k$, let $A_k=\{\pi^{-1}(k):\pi\in P\}$. This set
is nonempty. On ambiguous inputs, membership of $k$ is either constantly
true, or the test $S\cap A_k\ne\varnothing$. With three roles, if
$|A_k|\geq2$, that test is always true. Hence every nonconstant output
coordinate is exactly one input membership bit, say $k\in f(S)$ iff
$i\in S$. Record $p(i)=k$. The recorded sources are distinct, because
every permutation in $P$ must send each such source to its recorded
target. Thus $p$ is a partial permutation, and the other coordinates
form $C_p$.

The atomic upper bound above excludes each variable output coordinate
from all rows except its source row. The mandatory-image
fact~\eqref{sym:eq:mandatory} includes it in that row. These are precisely
the bounds in Theorem~\ref{sym:thm:unary}. The partial permutation is
uniquely determined by its action on the three pairs: a variable output
coordinate is absent exactly on the pair omitting its source. Atomic
rows are actual values, so they too are unique.

\subsubsection{Synthesis of every normal form}

Use the already established reduct construction
\[
 \Reg(C,E)(S)=\begin{cases}\{i\}\cup E_i,&S=\{i\},\\
 S\cup C,&|S|\geq2,\end{cases}\qquad E_i\subseteq C\setminus\{i\}.
\]
It uses only singleton constants and owned choice. Its generic syntactic
synthesis is proved in \texttt{Synthesis.lean}.

\emph{Rank three.} The partial permutation is total, $C_p$ is empty,
and every row is its prescribed singleton. Use that role permutation.

\emph{Rank two.} Extend $p$ to a total permutation $\pi$. Its remaining
output is the single element of $C_p$. Apply a reduct function $\Reg(C_p,E)$
after $\pi(x)$, choosing each row's optional extra winner as prescribed.
The remaining input row is the singleton in $C_p$. This realizes all
four row choices for each rank-two partial permutation.

\emph{Rank one.} By permuting the input, normalize $p$ to $p(a)=a$ and
$C_p=\{b,c\}$. First choose any desired nonempty rows $D_b,D_c$ within
$\{b,c\}$ while fixing the row $a$. On inputs $b,c$, owned choices by
$b,c$ are Boolean disjunction and conjunction, and the transposition
$(b\ c)$ is Boolean negation. Relative constants are
$x\ow b(b\ c)(x)$ and $x\ow c(b\ c)(x)$. Every expression so built
still sends $\{a\}$ to $\{a\}$. Taking unions of two Boolean functions
realizes any prescribed pair of nonempty rows $D_b,D_c$.

Precompose this expression with $\Reg(C_p,0)$, which fixes all atomic inputs
and sends ambiguous inputs to $C_p$ or $I$. The Boolean expressions
above fix both those sets, so the resulting ambiguous behavior is
exactly $F_p$. Finally, to add any prescribed $Z\subseteq C_p$ to the
$a$ row, postcompose with $\Reg(Z,E)$ having $E_a=Z$ and the other rows
empty. It preserves every nonempty subset of $C_p$, as well as $I$.
Thus all four choices for the $a$ row and all nine choices for the other
two rows are realized independently.

\emph{Rank zero.} The ambiguous output is always $I$ and the three
atomic rows may be any nonempty sets. The term
\[
 x\ow a(a\ b)(x)
\]
sends singleton $b$ to $a$ and fixes the other singleton roles. It
preserves $I$. These elementary mergers, together with permutations,
realize every map $h:I\to I$: merge each input into a chosen
representative of its $h$-fiber, then permute the distinct representatives
to their desired images. Representatives are never merged, so the
construction has the asserted effect. Union at most three such maps
to realize arbitrary nonempty atomic rows, and precompose with $\Reg(I,0)$.
All ambiguous inputs then become $I$, while atomic rows are unchanged.

This proves synthesis. Larger ground constants give the constant
functions. Any nonconstant live function takes $I$ to $I$, so the only
constant overlap is $I$ itself.

\begin{corollary}[Exact count]
There are exactly $751$ unary term operations in the three-role full
signature. Of these, $745$ preserve $I$.
\end{corollary}
\begin{proof}
For rank $r$, the number of partial permutations is
$\binom3r^2r!$. The number of row choices for each is:
\[
\begin{array}{c|rrrr}
 r&0&1&2&3\\\hline
 \text{partial permutations}&1&9&18&6\\
 \text{row choices each}&343&36&4&1
\end{array}
\]
Hence there are $343+324+72+6=745$ normal forms of
type~\eqref{sym:eq:unary}. Add the six constant functions other than $I$.
\end{proof}

\subsection{Context composition and the symmetric inverse monoid}

The partial permutations of $I$, including the empty map, form the
standard symmetric inverse monoid $\mathcal I_I$ under partial
composition~\cite{East2006}. For three roles its size is
$\sum_{r=0}^3\binom3r^2r!=34$.

\begin{theorem}[An ambiguity-action quotient]\label{sym:thm:monoid}
The $745$ unary term operations preserving $I$ form a monoid under
composition. Sending the normal form $(p,D)$ to $p$ is a surjective
monoid homomorphism onto $\mathcal I_I$.
\end{theorem}
\begin{proof}
Every ambiguous input has an ambiguous output, and direct calculation
gives $F_p\circ F_q=F_{p\circ q}$. Indeed, an output coordinate either
is fixed present, or follows a unique input coordinate through each of
the two partial permutations. The atomic rows of a composite are
\[
 D^{f\circ g}_i=\begin{cases}
 D^f_j,&D^g_i=\{j\},\\
 F_{p_f}(D^g_i),&|D^g_i|\geq2.
 \end{cases}
\]
These formulas give the exact composite. Every $p$ has an admissible row
family and is realized by the synthesis, so the homomorphism is onto.
\end{proof}

\begin{corollary}[All reversible unary operations]
The invertible unary term operations are exactly the six role
permutations.
\end{corollary}
\begin{proof}
An invertible function fixes $I$: every other possibility is a constant.
Its partial-permutation image must be a unit, hence a total permutation.
Rank three leaves no choice of atomic rows. Conversely every role
permutation has its inverse as a term.
\end{proof}

Thus the familiar inverse monoid is a quotient governing ambiguity,
not the whole algebra of unary contexts. Its fibers retain the permitted
atomic rows. This is genuine context composition with internal
role-switching; it does not use uniform-profile selection, whose lift
failed to preserve substitution in \cref{app:contexts}.

\subsection{The known relation monoid and its affine extension}

Regard a relation $R\subseteq I\times I$ as sending source roles to
output roles, and write $R[S]$ for its direct image. Let $\FP(\Sym(I))$
be the monoid of nonempty unions of permutation graphs, with relational
composition. This is the \emph{factorpower} of the natural symmetric-group
action, introduced by Ganyushkin and Mazorchuk and surveyed by
Mazorchuk\cite{GanyushkinMazorchuk1995,Mazorchuk2009}. In the usual
definition, nonempty subsets $P,Q\subseteq\Sym(I)$ are identified when
$\{\pi(i):\pi\in P\}=\{\pi(i):\pi\in Q\}$ for every $i$.
Thus the relation monoid appearing here is exactly a known object.

Form its affine extension
\[
 \mathcal E_I=\mathcal P(I)\rtimes\FP(\Sym(I)),\qquad
 (C,R)(D,Q)=(C\cup R[D],R\circ Q).
\]
The action on $\mathcal P(I)$ is by direct image; it preserves empty set
and unions, so this multiplication is associative. Its identity is
$(\varnothing,\mathrm{id})$. The affine action is
\begin{equation}\label{gen:eq:affine}
 F_{C,R}(S)=C\cup R[S].
\end{equation}
Every such map sends ambiguous sets to ambiguous sets because $R$
contains a permutation. Restrict these maps to
$\mathcal A_I=\{S\subseteq I:|S|\geq2\}$ and identify equal restrictions.
Call the resulting transformation monoid $\mathcal Q_I$.

\begin{theorem}[Exact ambiguity quotient]\label{gen:thm:quotient}
Let $\mathcal M_I$ be the monoid of unary term functions fixing $I$.
Restriction to $\mathcal A_I$ is a surjective monoid homomorphism
$\mathcal M_I\twoheadrightarrow\mathcal Q_I$.
\end{theorem}
\begin{proof}
A unary term is either globally constant or has live provenance
$C\cup\bigcup_{\pi\in P}\pi S$ on ambiguous inputs, with $P$ nonempty.
Here is a simultaneous induction establishing that dichotomy. The
variable is live and a winner is constant. Switching transports constants
and permutations. At an owned node, a live child is always ambiguous on
ambiguous inputs. Thus collapse there occurs exactly when a globally
constant child is the owner singleton, making the whole node globally
constant. Otherwise union the live provenances and the constant children.
A node with only constant children is again globally constant.

Every live expression fixes $I$. Hence a top-preserving unary term is
live or is the constant $I$, and has a restriction of the displayed form.
Conversely, for any $(C,R)$ choose a nonempty finite set of permutations
whose union is $R$. The term joining their images of $x$ and the
singleton constants in $C$ realizes \eqref{gen:eq:affine} on all inputs and
fixes $I$. Restriction respects composition since all these maps preserve
$\mathcal A_I$.
\end{proof}

The constant part cannot be suppressed. For example,
$S\mapsto S\cup\{a\}$ on ambiguous inputs cannot equal $R[S]$ for a
union of permutations $R$. Every other output coordinate $b\ne a$ has
the singleton source $b$. As shown below, that source is uniquely
determined by the restricted action. Thus every constituent permutation
must fix all $b\ne a$, hence also $a$. Their union is the identity and
cannot make $a$ constantly present. The relevant object is an affine
extension and a quotient of a known monoid, not just the factorpower.

\subsection{Canonical rows, matching, and synthesis}

For an output role $k$ let
$R^{-1}(k)=\{i:(i,k)\in R\}$. Membership of $k$ in
\eqref{gen:eq:affine} is either constant or the predicate
$S\cap R^{-1}(k)\ne\varnothing$. Every such source set is nonempty.
On ambiguous inputs this predicate is constantly true exactly when its
source set has size at least $K-1$.

\begin{lemma}[Unique row representation]\label{gen:lem:rows}
Every action in $\mathcal Q_I$ has a unique row for each output role,
of one of the following forms:
\[
 *\quad\text{(always present)},\qquad
 A_k\subseteq I,\quad 1\leq|A_k|\leq K-2
 \quad\text{(present iff }S\cap A_k\ne\varnothing\text{)}.
\]
\end{lemma}
\begin{proof}
The preceding threshold gives existence. A nonconstant row has a pair
disjoint from its source set, distinguishing it from $*$. If two small
source sets $A,B$ differ, choose $i\in A\setminus B$, swapping their
names if necessary. Since $|B|\leq K-2$, choose a different
$j\notin B$. The pair $\{i,j\}$ meets $A$ and misses $B$, distinguishing
the rows.
\end{proof}

Not every syntactically possible collection of rows is realizable. Build
a bipartite graph $H$ with source and output copies of $I$. A nonconstant
output row $k$ allows exactly the edges from $A_k$; a constant row allows
edges from every source. Edges into nonconstant rows are
\emph{prescribed}; edges into constant rows are only allowed.

\begin{theorem}[Matching criterion]\label{gen:thm:matching}
A collection of canonical rows belongs to $\mathcal Q_I$ exactly when
every prescribed edge belongs to a perfect matching of $H$. If all rows
are constant, it always belongs.
\end{theorem}
\begin{proof}
For necessity, choose a representative $(C,R)$. A nonconstant row
cannot belong to $C$, and its preimage under $R$ is exactly its canonical
source set by Lemma~\ref{gen:lem:rows}. Every prescribed edge lies in some
permutation constituent of $R$. That permutation is a perfect matching
of $H$.

Conversely, choose a perfect matching for each prescribed edge. Their
union $R$ is a union of permutations. At a nonconstant row its preimage
is contained in the allowed source set and includes every prescribed
source, so it is exactly $A_k$. Set $C$ equal to the set of constant
output rows. Then $(C,R)$ realizes the given rows. If no edges are
prescribed, take $R$ to be the identity and $C=I$.
\end{proof}

For comparison, the analogous criterion with \emph{all} relation edges
prescribed is the standard factorpower realizability criterion
(Mazorchuk\cite[Section~2.1]{Mazorchuk2009}). Here the constant rows leave
their matching edges free, and the ambiguity threshold removes larger
nonconstant source sets. An ordinary Hall existence test alone is
insufficient. At $K=4$, rows
\[
 (\{1\},\{1,2\},*,*)
\]
allow a perfect matching, but the prescribed edge from source $1$ to
output $2$ belongs to none: output $1$ already requires source $1$.
The collection is therefore not realizable.

\begin{corollary}[Polynomial recognition and synthesis]
Given $K$ canonical rows, membership in $\mathcal Q_I$ is decidable in
polynomial time. Every member has a realization by at most $K^2$
permutation images, together with its constant roles, and by an ordinary
term of size polynomial in $K$.
\end{corollary}
\begin{proof}
For each of the at most $K(K-2)$ prescribed edges, force that edge and
find a perfect matching in the remaining graph. Standard augmenting
paths suffice. Use those witness permutations in the synthesis from
Theorem~\ref{gen:thm:quotient}. Join the pieces in a balanced tree. Expanding
the fixed derived-union term duplicates each argument twice; balanced
joining of $m$ pieces has size $O(m^2)$ rather than the exponential
expansion of a skew join. Singleton constant pieces and explicit
permutation descriptions have polynomial size too.
\end{proof}

\begin{theorem}[Canonical composition]\label{gen:thm:composition}
For outer rows $A$ and inner rows $B$, the composite row at $k$ is:
\begin{enumerate}
\item $*$ if $A_k=*$ or some $i\in A_k$ has $B_i=*$;
\item otherwise $\bigcup_{i\in A_k}B_i$, changed to $*$ if that union
has at least $K-1$ elements.
\end{enumerate}
This computes composition in time polynomial in $K$ on descriptions of
$O(K^2)$ bits.
\end{theorem}
\begin{proof}
Substitute the inner membership disjunctions into the outer membership
disjunction. A constant true disjunct makes the result true; otherwise
the sources union. Apply the ambiguity threshold. Closure follows also
from the affine multiplication in $\mathcal E_I$.
\end{proof}

Together, \cref{gen:lem:rows,gen:thm:matching,gen:thm:composition} give
an explicit model of the entire ambiguity-action monoid $\mathcal Q_I$:
its elements are exactly the admissible canonical row tables, its
identity has row $\{k\}$ at output $k$, and its multiplication is the
displayed row substitution and threshold rule. This is a canonical
representation with effective multiplication, not a presentation by
abstract generators and defining relations. It classifies the whole
ambiguity quotient for every $K$; what remains unclassified is the
general singleton fiber above an action, and hence the full unary
context monoid at larger $K$.

For three roles every nonconstant row is a singleton. The matching
criterion says that their sources must be distinct. Thus these rows
are exactly a partial bijection, with $*$ on the missing output roles,
recovering the $34$-element symmetric inverse monoid. Four roles admit
two-source rows, and the matching criterion is a genuine further
constraint. Exhausting the criterion and, independently, all affine
permutation-union representatives gives
\[
\begin{array}{c|rr}
K&|\FP(\Sym(I))|&|\mathcal Q_I|\\\hline
3&49&34\\
4&7443&1979
\end{array}
\]
These are finite computations, with complete enumerations reproduced by
\texttt{check\_general.py}. They count the ambiguity quotient, not the
full unary term monoid. Its singleton fibers for general $K$ are still
unclassified.

\subsection{Relative reset and exact two-permutation fibers}

\subsubsection{A uniform reset operation}

Given $A,B$, start with $A\cup B$ and, once for each owner $r$, replace
the current value $U$ by $U\star_r B$. Write the resulting derived
binary operation as $\Delta(A,B)$.

\begin{lemma}[Relative reset]\label{gen:lem:relative}
\[
 \Delta(A,B)=
 \begin{cases}B,&|B|=1,\\A\cup B,&|B|\geq2.\end{cases}
\]
Its value is independent of the enumeration of owners. It is
idempotent and equivariant.
\end{lemma}
\begin{proof}
Every intermediate value contains $B$. If $B$ is ambiguous, both
operands are ambiguous, and each step preserves $A\cup B$. If
$B=\{r\}$, its own step forces that singleton; every later step leaves
it fixed. Idempotency and equivariance follow from the displayed law.
The atomic and ambiguous cases are proved in \texttt{Separation.lean}.
\end{proof}

This construction also gives a precise observational quotient.
For an idempotent atom-preserving term, record its support $\Lambda$ on
ambiguous inputs and its projected argument $j$ on atomic inputs.
The extracted core cannot depend on an absent argument, so $j\in \Lambda$.

\begin{proposition}[Homogeneous observation quotient]\label{gen:prop:pointed}
These pairs are exactly $(\Lambda,j)$ with $\varnothing\ne \Lambda$ and $j\in \Lambda$.
They compose by
\[
 (\Lambda,j)\bigl((L_i,\ell_i)_i\bigr)
   =\left(\bigcup_{i\in \Lambda}L_i,\ell_j\right).
\]
They form the term clone of right normal bands. Relative reset maps to
its binary multiplication.
\end{proposition}
\begin{proof}
Atomic and all-ambiguous input regions are invariant under these
operations, so their observations respect substitution. Every pair is
realized by $\Delta(\bigcup_{i\in \Lambda}x_i,x_j)$, and distinct pairs are
separated by atomic valuations or by a pair and $I$ on ambiguous inputs.
The binary operation on pairs is $(\Lambda,j)(L,\ell)=(\Lambda\cup L,\ell)$:
word terms remember their variable set and their last variable.
Associativity, $xx=x$, and $xyz=yxz$ present exactly these word
invariants, the standard right-normal-band
calculus\cite[Section~8]{RomanowskaSmith1989}.
\end{proof}

In the notation of Freese et al., this is the term clone of
$\mathcal P_{\vee}$, the unique largest proper subvariety of the
reset variety $\mathcal P_{2+1}$
\cite[Theorems~3.1, 3.7--3.8]{FreeseEtAl2002}.

This is a quotient of operations observed on two invariant regions,
not a quotient identifying winner-set values. Mixed valuations retain
more information. For distinct $a,b,c$,
\[
 \Delta(\Delta(\{c\},\{a\}),\{a,b\})=\{a,b\},\qquad
 \Delta(\{c\},\Delta(\{a\},\{a,b\}))=\{a,b,c\}.
\]
Thus $\Delta$ is not associative in the full algebra. The right-normal-band
presentation is complete for the stated observation quotient only.
\Cref{sec:reset} gives the complete mixed-input
semantics of the reset-only fragment: antimatroids with a unique initial
element, polynomial prerequisite normal forms, and a three-element
generator identified with Waldhauser's $G_3^{\rm op}$\cite{Waldhauser2006}.

\subsubsection{Marginal atomic bounds do not suffice}

For $K=4$ let $\sigma=(1\ 2\ 3\ 4)$ and fix the ambiguity action
$F(S)=S\cup\sigma S$. Each atomic row must be a nonempty subset of
$\{i,\sigma(i)\}$, with no mandatory singleton-source image. These
obvious marginal bounds allow the atomic function
\[
 h(1)=1,\quad h(2)=2,\quad h(3)=3,\quad h(4)=1.
\]
Nevertheless no term has these rows and ambiguity action $F$.
We prove a uniform classification explaining the obstruction.

\begin{lemma}[Localization without constant coordinates]\label{gen:lem:local}
Suppose a unary term has ambiguity action $R[S]$, with no constant
output coordinate. It contains no winner leaves. Push switches to its
variable leaves.
Every remaining leaf permutation lies in the perfect matchings of the
canonical support graph.
\end{lemma}
\begin{proof}
On ambiguous inputs a live term has constant-plus-permutation provenance.
Any constant contribution would make an output coordinate constant,
so it is empty. As in Theorem~\ref{gen:thm:core}, any winner leaf would force
a nonempty constant contribution, so there are none. Each leaf
permutation contributes its graph to $R$. For nonconstant rows the
canonical source sets are uniquely determined, so its graph lies in
the allowed graph. Conversely every such permutation is a perfect
matching and an allowed generator.
\end{proof}

For $R=\mathrm{id}\cup\sigma$ and $K\geq4$, every source row has size
one or two, hence is nonconstant. Its perfect matchings independently
choose identity or $\sigma$ on each nontrivial cycle: choosing a shifted
edge forces the next shifted edge around that cycle.

For a cycle $C$, describe an atomic row at $i\in C$ by
$0$ for $\{i\}$, $1$ for $\{\sigma(i)\}$, and $?$ for their pair.
A fixed point has just its one mandatory singleton.

\begin{theorem}[Exact two-permutation fibers]\label{gen:thm:cycles}
Fix $K\geq4$ and the ambiguity action $S\mapsto S\cup\sigma S$.
The choices on different cycles of $\sigma$ are independent.
On a cycle of length $m$ the permitted atomic tables are exactly:
\begin{itemize}
\item for $m=1$, its fixed singleton;
\item for $m=2$, every word in $\{0,1,?\}^2$;
\item for $m\geq3$, every word in $\{0,?\}^m$ or $\{1,?\}^m$,
and the $m$ words with $0$ at $r$, $1$ at $\sigma^{-1}(r)$,
and $?$ elsewhere.
\end{itemize}
Consequently the fiber size is
\[
 \prod_{C\text{ a cycle of }\sigma}g(|C|),\qquad
 g(1)=1,\quad g(2)=9,\quad g(m)=2^{m+1}+m-1\ (m\geq3).
\]
An arbitrary action $p[S]\cup q[S]$ is covered by normalizing the output
with $p^{-1}$ and taking $\sigma=p^{-1}q$.
\end{theorem}
\begin{proof}
For necessity, the atomic upper bound gives the displayed row options.
Suppose a long cycle has row $0$ at $i$ and row $1$ at $j$. If
$i\ne\sigma(j)$, these are distinct singleton winners. The common-source
theorem requires one allowed leaf permutation producing both. But an
allowed matching on this cycle is either identity everywhere or shifted
everywhere. Neither does so. Thus $i=\sigma(j)$. If both kinds of
singleton occur there can be only one of each, giving precisely the
listed exceptional words. For a fixed point nonemptiness forces its
singleton; for a two-cycle this argument imposes no further restriction.

For synthesis on one cycle, start with the union of its two permutation
generators, giving all $?$ rows. To force $0$ at selected positions,
successively apply $U\mapsto U\star_i x$. This affects only the chosen
row and retains earlier $0$ rows. To force $1$ instead, use
$U\mapsto U\star_{\sigma(i)}\sigma(x)$. The exceptional word at $r$ is
$x\star_r\sigma(x)$. These constructions also realize all nine words
on a two-cycle.

The allowed leaf permutations already form the Cartesian product of
the two orientations on each cycle. All owned operations and union are
idempotent. Closure of a Cartesian product under idempotent
coordinatewise operations is the product of the local closures.
For two coordinates, suppose $t$ realizes $a$ from generators
$a_1,\ldots,a_n$ and $u$ realizes $b$ from $b_1,\ldots,b_m$.
Choose $z_{ij}=(a_i,b_j)$ from the full Cartesian product of generators.
Then
\[
 t\bigl(u(z_{11},\ldots,z_{1m}),\ldots,
         u(z_{n1},\ldots,z_{nm})\bigr)=(a,b):
\]
each inner application preserves $a_i$ by idempotency, and the outer
application preserves $b$ by idempotency. Induct on the number of
coordinates. Thus these local atomic tables
can be realized independently, by a term $t$ whose ambiguity support
lies in $\mathrm{id}\cup\sigma$.

Finally apply $\Delta(E,t)$ where $E(x)=x\cup\sigma(x)$.
On ambiguous inputs $t(S)\subseteq E(S)$ and $t(S)$ is ambiguous, so
reset retains exactly $E(S)$. On an atomic input the prescribed row
$t(\{i\})$ is either a singleton or the full allowed pair. In the first
case reset selects it; in the second $E(\{i\})$ already equals it.
This restores full ambiguity support without changing any atomic row.
The two uniform-orientation families overlap once; the $m$ exceptional
words are separate. This proves the count and independence.
\end{proof}

The example above violates this coherence condition at inputs $2$ and
$4$. A four-cycle has $35$ permitted tables, not the $3^4=81$ allowed
by its marginal bounds. Longer cycles give $68$ and $133$ at lengths
five and six. If every atomic row is a singleton, a long cycle must
choose one orientation everywhere; a two-cycle may also merge its two
roles.

\begin{corollary}[All four-role fibers without constant coordinates]
At $K=4$, every ambiguity action without constant coordinates is a union
of two permutations. There are $282$ such actions and exactly $6378$
unary term functions above them.
\end{corollary}
\begin{proof}
Every canonical row has at most two sources. A singleton row forces its
matched source to occur nowhere else, since every edge extends to a
matching. Remove such pairs. In the remaining support graph all output
degrees are two, while every input degree is at least two: an input of
degree one would force its output to have degree one. Counting edges
makes all input degrees two as well. Components are even cycles, each
the union of two matchings. Apply Theorem~\ref{gen:thm:cycles}.

For normalized cycle types $1^4,\,2\,1^2,\,2^2,\,3\,1,\,4$, the numbers
of distinct support graphs are respectively $24,72,18,96,72$. These
follow by choosing the first permutation and the cycle permutation, then
dividing by the two orientation choices in each nontrivial component.
They sum to $282$. The fiber sizes are respectively $1,9,81,18,35$,
giving $6378$ in total. The enumeration in \texttt{check\_fibers.py}
independently checks every four-role support.
\end{proof}

\subsection{Reversible contexts at every number of roles}

\begin{lemma}
Every term function is monotone for inclusion on nonempty winner sets.
\end{lemma}
\begin{proof}
Constants, variables, and permutation images are monotone. Suppose
$S\subseteq S'$ and $T\subseteq T'$, all nonempty. If an upper operand
equals $\{r\}$, its corresponding lower operand is also $\{r\}$, and
both owned outputs equal $\{r\}$. Otherwise the upper output is
$S'\cup T'$. The lower output is either $S\cup T$, which is contained
in it, or $\{r\}$ supplied by one lower operand, which is contained in
the corresponding upper operand. Thus owned choice is monotone.
\end{proof}

\begin{theorem}[All unary units]
For every $K\geq3$, the group of units of the full unary term monoid is
exactly $\Sym(I)$.
\end{theorem}
\begin{proof}
A unit and its term inverse are monotone, so they form an order
automorphism of the poset $\Pow(I)$. Such an automorphism permutes its
minimal elements, the singletons. Every set is determined by the
singletons below it, so that permutation determines the whole map.
Conversely every role permutation and its inverse are terms.
\end{proof}

This generalizes the three-role unit calculation without requiring a
higher-role classification of singleton fibers. It is a direct use of
the familiar order-automorphism argument for the Boolean subset poset.
